\chapter{LİTERATÜR TARAMASI}
\thispagestyle{empty}

Bu bölümde, projenin dayandığı temel kavramlar ve ilgili çalışmalar tematik olarak incelenmiştir. İnceleme dört ana eksende düzenlenmiştir: graf sinir ağı mimarileri, metin temsili ve dil modelleri, çok kipli öneri sistemleri ve Wikipedia üzerine uygulamalı çalışmalar. Bölümün sonunda, incelenen on iki çalışmanın karşılaştırmalı bir özeti tablo halinde sunulmaktadır.

\section{Graf Sinir Ağı Mimarileri}

Graf yapılı veriler üzerinde derin öğrenme yöntemlerinin uygulanması, son yıllarda hızla gelişen bir araştırma alanıdır. Bu alandaki temel mimariler, düğümler arasındaki bilgi yayılımını farklı stratejilerle gerçekleştirmektedir.

Kipf ve Welling \cite{kipf2017semi} tarafından önerilen \textbf{Graph Convolutional Networks (GCN)}, spektral graf konvolüsyonlarının birinci dereceden yaklaşıklığını kullanarak graflar üzerinde yarı denetimli (semi-supervised) öğrenmeyi verimli hale getirmiştir. GCN'in yayılım kuralı
\[
    H^{(l+1)} = \sigma\!\left(\tilde{D}^{-\frac{1}{2}}\tilde{A}\,\tilde{D}^{-\frac{1}{2}}H^{(l)}W^{(l)}\right)
\]
biçiminde tanımlanır; burada $\tilde{A}$, self-loop eklenmiş komşuluk matrisidir. Model, Cora, Citeseer ve Pubmed atıf ağlarında önceki yarı denetimli yöntemleri geride bırakmış ve kenar sayısına göre doğrusal ($\mathcal{O}(|\mathcal{E}|)$) hesaplama karmaşıklığı sayesinde büyük graflarda hızlı eğitim olanağı sunmuştur. Bununla birlikte GCN, her komşuya düğümün derecesine göre normalize edilmiş aynı ağırlığı atfetmekte ve yönlü grafları doğrudan desteklememektedir.

Bu sınırlamayı gidermek amacıyla Veličković vd.\@ \cite{velickovic2018graph} \textbf{Graph Attention Networks (GAT)} modelini önermiştir. GAT, Transformer mimarisindeki öz-dikkat (self-attention) mekanizmasını graf yapısına uyarlayarak her komşu düğüme farklı dikkat ağırlıkları atamaktadır. Çok başlıklı dikkat (multi-head attention) yapısı, modelin farklı alt uzaylardaki ilişkileri paralel olarak öğrenmesini sağlar. GAT, transdüktif atıf ağı görevlerinde GCN'e göre \%1,5--2 oranında iyileştirme sağlamış; indüktif PPI veri setinde ise GraphSAGE'e kıyasla \%20,5'lik bir performans artışı sergilemiştir. Dikkat ağırlıklarının görselleştirilebilir olması, modelin hangi bağlantılara odaklandığını yorumlanabilir kılmaktadır. Wikipedia bağlamında bu özellik kritik önem taşır: bir makaledeki her hiperlink içeriğin özüyle aynı derecede alakalı değildir ve GAT bu farkı otomatik olarak öğrenebilmektedir.

Hamilton, Ying ve Leskovec \cite{hamilton2017inductive} tarafından geliştirilen \textbf{GraphSAGE} ise farklı bir tasarım felsefesi sunmaktadır. Model, tüm komşuları işlemek yerine sabit sayıda komşu örnekleyerek (sampling) bellek kullanımını kontrol altında tutar ve eğitilebilir birleştirme fonksiyonları (mean, LSTM, pooling aggregator) aracılığıyla düğüm temsillerini oluşturur. En önemli özelliği indüktif yapısıdır: eğitim sırasında hiç görülmemiş düğümler için de embedding üretebilmektedir. GraphSAGE, DeepWalk'a kıyasla F1 skorlarında ortalama \%51 artış sağlamıştır. Ancak rastgele komşu örneklemesi kritik komşuların göz ardı edilmesine yol açabilmekte ve performans artışı genellikle iki hop derinliğinden sonra durmaktadır.

Bu projede birincil model olarak \textbf{GAT} tercih edilmiştir. Bunun başlıca nedeni, Wikipedia hiperlink grafındaki bağlantıların anlamsal ağırlıklarının farklılık göstermesi ve dikkat mekanizmasının bu farklılığı doğrudan modelleyebilmesidir. Karşılaştırma modeli (baseline) olarak ise GCN kullanılacaktır.

\section{Metin Temsili ve Dil Modelleri}

Graf sinir ağlarının düğüm öznitelikleri olarak kullanılacak metin temsillerinin kalitesi, öneri sisteminin başarısını doğrudan etkilemektedir. Bu alandaki gelişmeler, Transformer mimarisi ile başlayan bir evrim sürecinin ürünüdür.

Vaswani vd.\@ \cite{vaswani2017attention} tarafından önerilen \textbf{Transformer} mimarisi, yinelemeli (RNN) ve evrişimli (CNN) yapıları tamamen terk ederek yalnızca dikkat mekanizmalarına dayanan bir kodlayıcı-kod çözücü yapısı sunmuştur. Öz-dikkat (self-attention) katmanı,
\[
    \text{Attention}(Q, K, V) = \text{softmax}\!\left(\frac{QK^T}{\sqrt{d_k}}\right)V
\]
formülü ile dizideki her konumun diğer tüm konumlarla ilişkisini tek adımda hesaplar. Bu yapı, uzun mesafeli bağımlılıkların yakalanmasını kolaylaştırmış ve paralel hesaplama sayesinde eğitim süresini önemli ölçüde kısaltmıştır. WMT 2014 çeviri görevinde Transformer, önceki en iyi modellerin eğitim maliyetinin dörtte biriyle yeni performans rekorları kırmıştır.

Devlin vd.\@ \cite{devlin2019bert}, Transformer'ın yalnızca kodlayıcı (encoder) katmanlarını kullanarak \textbf{BERT} modelini geliştirmiştir. Maskelenmiş Dil Modeli (Masked Language Model) ve Sonraki Cümle Tahmini (Next Sentence Prediction) görevleri ile çift yönlü ön eğitim yapan BERT, yayınlandığı dönemde 11 farklı NLP görevinde en iyi sonuçları (SOTA) elde etmiştir. BERT'in bağlamsal temsil gücü -- aynı kelimenin farklı cümlelerdeki anlamını ayırt edebilmesi -- metin gömme alanında bir dönüm noktası olmuştur. Ancak 512 token girdi sınırı ve yüksek hesaplama maliyeti, büyük ölçekli uygulamalarda doğrudan kullanımını zorlaştırmaktadır.

BERT'in ardından gelen gelişmeler, cümle düzeyinde anlamsal benzerlik ve yoğun bilgi erişimi (dense retrieval) için özelleştirilmiş \textbf{cümle gömme modellerini} ortaya çıkarmıştır. Bu projede iki modern cümle gömme modeli karşılaştırmalı olarak değerlendirilmektedir:
\begin{itemize}
    \item \textbf{BAAI/bge-m3} \cite{chen2025bgem3}: Çok dillilik, çok işlevsellik ve çok ayrıntılılık (multi-linguality, multi-functionality, multi-granularity) ilkelerine göre tasarlanmış bir cümle gömme modelidir. Öz-bilgi damıtma (self-knowledge distillation) yöntemiyle eğitilen model, yoğun (dense), seyrek (sparse) ve çoklu vektör (multi-vector) erişim yöntemlerini tek bir çerçevede birleştirmektedir. 100'den fazla dili destekleyen bge-m3, MIRACL ve MKQA gibi çok dilli erişim kıyaslamalarında önceki en iyi sonuçları geride bırakmış ve 1024 boyutlu gömme vektörleri üretmektedir.
    \item \textbf{Alibaba-NLP/gte-modernbert-base} \cite{zhang2024mgte}: GTE (General Text Embeddings) ailesi kapsamında geliştirilen mGTE, RoPE konum kodlaması ve unpadding optimizasyonu ile güçlendirilmiş bir metin kodlayıcı üzerine inşa edilmiştir. Model, 8192 token uzunluğunda doğal bağlam penceresiyle ön eğitim görmüş olup uzun belge erişiminde aynı boyuttaki önceki en iyi model XLM-R'yi geride bırakmıştır. Projede kullanılan ModernBERT tabanlı varyant, 768 boyutlu gömme vektörleri üretmektedir.
\end{itemize}
Bu modeller, Transformer ve BERT'in bağlamsal temsil gücünü miras alırken, cümle düzeyinde anlam yakalama ve büyük ölçekli vektör karşılaştırması için optimize edilmiştir. Her iki model de Wikipedia makale metinlerini yoğun vektörlere dönüştürerek GNN'e düğüm özniteliği olarak sağlamaktadır.

\section{Çok Kipli Öneri Sistemleri ve GNN ile Metin İşleme}

Metin verisi ile graf yapısını birlikte kullanan çok kipli (multimodal) sistemlerin tasarımı, hem graf inşa stratejisi hem de modalite birleştirme (fusion) yöntemi açısından kritik kararlar gerektirmektedir. Bu alandaki iki kapsamlı literatür taraması, projenin mimari kararlarına doğrudan rehberlik etmiştir.

Wang, Ding ve Han \cite{li2022graph} tarafından hazırlanan GNN tabanlı metin sınıflandırma taraması, modelleri graf oluşturma mekanizmalarına göre iki ana kategoriye ayırmaktadır: \textit{derlem seviyesi} (corpus-level) ve \textit{belge seviyesi} (document-level). Derlem seviyesi yaklaşımlar (TextGCN, BertGCN gibi) tüm veri setini tek bir büyük graf olarak modeller ve kelime-belge ilişkilerini (TF-IDF, PMI) kenar olarak kullanır. Belge seviyesi yaklaşımlar ise her belge için ayrı bir graf oluşturur. Tarama, harici ön eğitimli modellerin (BERT gibi) GNN ile birleştirildiği mimarilerin en yüksek başarıyı sağladığını; ancak derlem seviyesi modellerin büyük veri setlerinde bellek darboğazına yol açtığını ortaya koymaktadır. Bu proje bağlamında Wikipedia'nın doğal hiperlink yapısı, harici bir graf inşasına gerek kalmadan doğrudan kullanılabilmekte ve ön eğitimli cümle gömme modelleri ile birleştirilerek hibrit bir yaklaşım oluşturulmaktadır.

Liu, Hu, Xiao vd.\@ \cite{zhou2023multimodal} tarafından sunulan çok kipli öneri sistemleri taraması ise bu alandaki modelleri üç aşamalı bir çerçevede sınıflandırmaktadır: \textit{modalite kodlama} (encoder), \textit{öznitelik etkileşimi} (interaction) ve \textit{öznitelik geliştirme} (enhancement). Öznitelik etkileşimi aşamasında graf yapılarının (kullanıcı-öğe ve öğe-öğe grafları) kullanımının, basit vektör birleştirmelerine (concatenation) göre çok daha etkili olduğu vurgulanmıştır. Ayrıca uçtan uca eğitimde büyük kodlayıcıların (BERT gibi) parametre maliyetinin yüksek olması nedeniyle iki aşamalı eğitim veya parametre dondurma stratejilerinin yaygın olarak tercih edildiği belirtilmiştir. Bu projede de benzer bir strateji izlenmektedir: cümle gömme modelleri önceden çalıştırılarak vektörler elde edilmekte ve GNN eğitimi bu sabit vektörler üzerinden gerçekleştirilmektedir.

\section{Wikipedia Üzerine Uygulamalı Çalışmalar}

Projenin doğrudan ilişkili olduğu üç uygulama çalışması, veri seti inşası, mimari tasarım ve metin-graf entegrasyonu konularında somut referanslar sunmaktadır.

Mernyei ve Cangea \cite{mernyei2020wikics} tarafından oluşturulan \textbf{Wiki-CS}, bilgisayar bilimi alanındaki 11.701 Wikipedia makalesini ve 216.123 hiperlinki içeren bir GNN kıyaslama veri setidir. Veri seti, Wikipedia'nın gürültülü kategori yapısının titizlikle temizlenmesiyle oluşturulmuş olup standart atıf ağlarına (Cora, Citeseer) göre çok daha yüksek bağlantı yoğunluğuna (ortalama derece: 36,94) sahiptir. Düğüm öznitelikleri olarak yalnızca GloVe kelime gömme ortalamalarının (300 boyutlu) kullanılması, metin temsilinin modern standartların gerisinde kalmasına neden olmaktadır. Bu proje, Wiki-CS'nin graf topolojisini temel alarak ancak metin temsilini modern cümle gömme modelleri ile yeniden oluşturarak \textbf{Custom-WikiCS} veri setini inşa etmiştir.

Moskalenko, Parra ve Saez-Trumper \cite{wikirecnet2020} tarafından geliştirilen \textbf{WikiRecNet}, 5,2 milyon düğüm ve 458 milyon kenar içeren tam ölçekli İngilizce Wikipedia üzerinde çalışan bir öneri sistemidir. Metin temsili için Doc2Vec, yapısal öğrenme için GraphSAGE ve hızlı aday erişimi için FAISS kullanılmıştır. Yapısal bilginin dahil edilmesi, yalnızca metin tabanlı önerilere göre nDCG@100 metriğinde yaklaşık \%100 iyileştirme sağlamıştır. Bu sonuç, metin ve graf bilgisinin birleştirilmesinin öneri kalitesini önemli ölçüde artırdığını doğrulamaktadır. WikiRecNet'in temel kısıtı, metin kodlayıcı olarak Doc2Vec gibi bağlam duyarsız bir yöntem kullanmasıdır; yazarlar bunu açıkça ``gelecek çalışma'' olarak not etmiştir. Bu proje, WikiRecNet'in genel mimarisinden ilham alırken Doc2Vec'i modern cümle gömme modelleri ile değiştirmektedir.

Javaji ve Sarode \cite{anime_hybrid_gnn_bert} tarafından sunulan hibrit öneri sistemi, Sentence-BERT gömmeleri ile GraphSAGE'i bir anime veri seti üzerinde birleştiren pratik bir uygulamadır. Çalışma, metin ve graf modalitelerinin kod düzeyinde nasıl entegre edileceğini (SequenceEncoder + GNNEncoder) somut olarak göstermektedir. Ancak test setinde yalnızca \%37 doğruluk elde edilmiş olması, hiperparametre optimizasyonu ve düzenlileştirme (regularization) eksikliğine işaret etmekte ve bu entegrasyonun dikkatli mühendislik gerektirdiğini ortaya koymaktadır.

\vspace{0.5em}
\noindent\textbf{Literatürün Genel Değerlendirmesi:}
İncelenen çalışmalar bir bütün olarak değerlendirildiğinde, literatürdeki üç temel boşluk belirginleşmektedir: (1)~Wikipedia üzerine yapılan en kapsamlı çalışma olan WikiRecNet'in bağlam duyarsız bir metin kodlayıcı (Doc2Vec) kullanması, (2)~GCN gibi temel mimarilerin bağlantı ağırlıklandırma yapamaması ve (3)~mevcut öneri sistemlerinin sonuçları düz bir liste olarak sunup kavramsal hiyerarşiyi göz ardı etmesi. Bu proje, modern cümle gömme modelleri ile GAT'ın dikkat mekanizmasını birleştirerek ilk iki boşluğu, LLM destekli hiyerarşik sunum katmanı ile üçüncü boşluğu gidermeyi hedeflemektedir. Tablo~\ref{tab:lit-karsilastirma}, incelenen çalışmaların karşılaştırmalı bir özetini sunmaktadır.

\begin{landscape}
\footnotesize
\setlength\tabcolsep{3pt}
\renewcommand{\arraystretch}{1.4}

\begin{longtable}{|p{2.0cm}|p{2.0cm}|p{2.8cm}|p{3.5cm}|p{3.0cm}|p{3.0cm}|p{3.8cm}|}

\caption{İncelenen çalışmaların karşılaştırmalı analizi.} \label{tab:lit-karsilastirma} \\
\hline
\textbf{Çalışma} & \textbf{Veri} & \textbf{Yöntem} & \textbf{Temel Bulgular} & \textbf{Güçlü Yönler} & \textbf{Kısıtlar} & \textbf{Proje ile İlişkisi} \\
\hline
\endfirsthead

\multicolumn{7}{c}{{\bfseries Tablo \thetable\ -- devamı}} \\
\hline
\textbf{Çalışma} & \textbf{Veri} & \textbf{Yöntem} & \textbf{Temel Bulgular} & \textbf{Güçlü Yönler} & \textbf{Kısıtlar} & \textbf{Proje ile İlişkisi} \\
\hline
\endhead

\hline \multicolumn{7}{r}{{Devamı sonraki sayfada}} \\ \hline
\endfoot

\hline
\endlastfoot

GAT \cite{velickovic2018graph} &
Cora, Citeseer, PPI &
Masked self-attention, multi-head &
GCN'e göre \%1,5--2 artış; PPI'da \%20,5 fark &
Komşu ağırlıklandırma, yorumlanabilirlik &
Head sayısı ile artan bellek maliyeti &
\textit{Birincil GNN modeli}; hiperlink ağırlıklandırması için kullanılacak \\
\hline

GCN \cite{kipf2017semi} &
Cora, Citeseer, Pubmed &
Spektral konv.\ 1.\ derece yaklaşım &
SOTA yarı denetimli sonuçlar; lineer ölçekleme &
Verimli, uçtan uca eğitilebilir &
Eşit ağırlık, yönlü graf desteği yok &
\textit{Baseline model} olarak karşılaştırma ölçütü \\
\hline

GraphSAGE \cite{hamilton2017inductive} &
Citation, Reddit, PPI &
Komşu örnekleme ve birleştirme &
F1'de ort.\ \%51 artış; indüktif başarı &
İndüktif öğrenme, ölçeklenebilirlik &
Örnekleme varyansı; 2-hop sonrası kazanç düşük &
\textit{İlişkili çalışma}; ölçeklenebilir GNN tasarımına referans \\
\hline

Transformer \cite{vaswani2017attention} &
WMT 2014 &
Self-attention, encoder-decoder &
1/4 maliyetle SOTA çeviri sonuçları &
Paralel eğitim, uzun mesafe yakalama &
$O(n^2)$ karmaşıklık &
\textit{Teorik temel}; GAT ve BERT'in dayandığı mimari \\
\hline

BERT \cite{devlin2019bert} &
Wikipedia + BooksCorpus &
MLM, NSP; çift yönlü ön eğitim &
11 NLP görevinde SOTA &
Bağlamsal temsil, alan uyumu &
512 token sınırı, yüksek hesaplama maliyeti &
\textit{Kavramsal öncül}; cümle gömme modellerinin temelini oluşturur \\
\hline

bge-m3 \cite{chen2025bgem3} &
MIRACL, MKQA, çok dilli &
Self-knowledge distillation; dense+sparse+multi-vec &
Çok dilli erişimde SOTA; 100+ dil desteği &
Tek çerçevede üç erişim yöntemi &
Yüksek boyut (1024), büyük model &
\textit{Birincil gömme modeli}; makale metinlerini vektörleştirir \\
\hline

mGTE \cite{zhang2024mgte} &
MTEB, çok dilli erişim &
RoPE + unpadding; 8192 token bağlam &
XLM-R'yi geride bırakır; uzun belgede avantaj &
Uzun bağlam, yüksek verimlilik &
Daha küçük boyut (768) &
\textit{İkincil gömme modeli}; karşılaştırma için kullanılır \\
\hline

GNN Metin Taraması \cite{li2022graph} &
20NG, MR, R8, SST-2 &
Corpus-level vs.\ document-level GNN &
BERT+GNN hibritleri en başarılı &
Kapsamlı taksonomi, graf inşa analizi &
Corpus-level bellek darboğazı &
\textit{Mimari rehber}; corpus-level hibrit yaklaşımı destekler \\
\hline

Multimodal RecSys Taraması \cite{zhou2023multimodal} &
Amazon, Tiktok, MovieLens &
Encoder-interaction-fusion çerçevesi &
Graf tabanlı etkileşim en etkili &
Cold-start çözümü, modern teknikler &
Gerçek zamanlı hesaplama karmaşıklığı &
\textit{Yol haritası}; çok kipli füzyon stratejisine temel \\
\hline

Wiki-CS \cite{mernyei2020wikics} &
Wikipedia CS (11.701 düğüm) &
Veri seti oluşturma, düğüm sınıfl. &
Yoğun yapı (ort.\ derece 36,94) &
Gerçekçi Wikipedia yapısı, 20 split &
GloVe ortalama temsili yetersiz &
\textit{Veri seti tabanı}; Custom-WikiCS'nin graf topoloji kaynağı \\
\hline

WikiRecNet \cite{wikirecnet2020} &
Wikipedia (5,2M düğüm) &
Doc2Vec + GraphSAGE + FAISS &
Yapısal bilgi ile \%100 nDCG artışı &
Ölçeklenebilirlik, indüktif yapı &
Eski metin kodlayıcı (Doc2Vec) &
\textit{Mimari taslak}; Doc2Vec $\to$ cümle gömme modernizasyonu \\
\hline

Hibrit Öneri \cite{anime_hybrid_gnn_bert} &
Anime Rec.\ DB (320K kullanıcı) &
Sentence-BERT + GraphSAGE &
Eğitimde \%52, testte \%37 doğruluk &
Metin+graf entegrasyon örneği &
Düşük genelleme, sınırlı metrikler &
\textit{Uygulama örneği}; entegrasyon derslerini sunar \\
\hline

\end{longtable}
\end{landscape}